% !TeX root = tesis.tex

Los pairings bilineales son una primitiva criptográfica fundamental en las construcciones modernas. Informalmente, un pairing es una aplicación que relaciona puntos de curvas elípticas con elementos de un grupo multiplicativo, preservando una estructura algebraica específica. Gracias a sus propiedades, los pairings permiten reducir ciertos problemas difíciles en curvas elípticas \cite{Menezes1993} (como el \ECDLP) a problemas equivalentes en otros grupos, lo que habilita nuevas construcciones criptográficas que no son posibles en el modelo clásico sin pairings.

\begin{definition}[Pairing bilineal]
	Sea $\E(\Fq)$ una curva elíptica definida sobre un cuerpo finito $\Fq$, y sea $r \in \N$ primo tal que $r$ divide a $|\E(\Fq)|$. Un pairing bilineal es una función
	\[
	e : \GOne \times \GTwo \to \GT,
	\]
	donde $\GOne,\GTwo$ son subgrupos cíclicos de orden $r$ de $\E(\Fq)$ (o de extensiones de $\E(\Fq)$), y $\GT$ es un subgrupo multiplicativo de un cuerpo finito (o de una extensión), también de orden $r$. Además, debe cumplir las siguientes propiedades:
	\begin{itemize}
		\item \textbf{Bilinealidad:} Para todos $P \in \GOne, Q \in \GTwo, a,b \in \Fr$, valen $\pairing{aP}{Q} = \pairing{P}{Q}^a$ y $\pairing{P}{bQ} = \pairing{P}{Q}^b$. Observar que en el lado derecho de la ecuación, se está utilizando la notación de grupo multiplicativo.
		\item \textbf{No degeneración:} Para cualesquiera generadores $G_1 \in \GOne, G_2 \in \GTwo$ de los subgrupos de orden $r$, $\pairing{G_1}{G_2} \neq 1_{\GT}$.
		\item \textbf{Computabilidad eficiente:} Existe un algoritmo eficiente en tiempo polinomial para computar $\pairing{P}{Q}$, para todos $P \in \GOne, Q \in \GTwo$.
	\end{itemize}
\end{definition}

Una forma útil de interpretar la bilinealidad es observar que el pairing bilineal convierte operaciones aditivas en operaciones multiplicativas:

$$\pairing{P + P'}{Q}  = \pairing{P}{Q} \cdot \pairing{P'}{Q}, \quad \pairing{P}{Q + Q'} = \pairing{P}{Q} \cdot \pairing{P}{Q'}.$$

Esto permite trasladar relaciones lineales entre puntos de curvas elípticas a relaciones multiplicativas en $\GT$, lo cual resulta extremadamente poderoso en construcciones criptográficas. Más aún, observamos que el pairing \textit{mezcla} los escalares $a,b$ en su producto, lo que induce otra operación difícil de revertir. Sin pairings, muchas construcciones modernas de pruebas sucintas y de autenticación algebraica no serían posibles.

La construcción explícita de un pairing bilineal involucra una cantidad considerable de teoría adicional, incluyendo el uso de cuerpos de extensión, divisores en curvas algebraicas, funciones racionales y algoritmos específicos como el algoritmo del loop de Miller \cite{miller86}. El desarrollo completo de estos conceptos excede el objetivo de esta tesis y requeriría una extensión significativa que no sería reutilizada en las secciones posteriores. Ejemplos clásicos de pairings bilineales utilizados en la práctica son el pairing de Weil y el pairing de Tate. Por lo tanto, se asumirá la existencia de los pairings bilineales junto con las propiedades mencionadas anteriormente, sin entrar en los detalles de su construcción matemática. Se remite al lector interesado a \cite{costello2010pairings}, que proporciona una introducción accesible a los pairings bilineales.

Los pairings bilineales que vamos a necesitar se dicen asimétricos, es decir, son tales que no se conoce ningún isomorfismo eficiente entre los subgrupos $\GOne,\GTwo$. Además, el \ECDLP debe ser difícil de resolver en ellos.
